\scanpage{17}
\mspage{\textcircled{3}}
\uline{NB} Si on \deleted{réduit}\inserted{étend} le groupe structural de
$Q_X$ sous $\mathcal N_{g,\nu}$ vers \inserted{son quotient}
$\Gamma_{\overline{\mathbb Q}/\mathbb Q}$, on trouve
\[
(18)\qquad R_X\overset{\mathrm{déf}}=
Q_X\wedge_{\mathcal N_{g,\nu}}\Gamma_{\overline{\mathbb Q}/\mathbb Q}
\simeq \unsure{R_X|Q_X}/\widehat\Gamma_{g,\nu}^{+}
\]
et ce \deleted{\unsure{Q}} $\Gamma_{\overline{\mathbb Q}/\mathbb Q}$-torseur
est can. isom. à
\[
(19)\qquad\left\{\begin{aligned}
R_X&\simeq\operatorname{Isom}_{\substack{\text{group. fond.}\deleted{\text{de}}\\\text{de Spec }\mathbb Q}}
(\operatorname{Spec}\overline{\mathbb Q},\operatorname{Spec}\widetilde{\mathbb Q})\\
&\simeq\operatorname{Isom}(\overline{\mathbb Q},\widetilde{\mathbb Q})\\
&\simeq\operatorname{plongts}(\overline{\mathbb Q},k).
\end{aligned}\right.
\]
Donc, la donnée d'une trivialisation de $R_X$, i.e. d'une restriction du
groupe structural de $Q_X$ de $\mathcal N_{g,\nu}$ à $\widehat\Gamma_{g,\nu}^{+}$
(ou de $P_{X,x}$ de $\mathcal M_{g,\nu}$ à $\widehat{\mathcal G}_{g,\nu}^{+}$)
équivaut à la donnée d'un iso.
\[
(20)\qquad\overline{\mathbb Q}\simeq\widetilde{\mathbb Q}
\qquad\text{i.e. plongmt }\overline{\mathbb Q}\subset k.
\]
\mspage{16}
